\lecture{7}
\title{Hash Functions (Part II)}

\begin{document}

Here's another way to handle key collisions.

\begin{solution}{Open Addressing}
    Every key $k \in \mathcal{U}$ is assigned not just one hash value in $\{0, 1, \dots, m - 1\}$, but rather a \emph{permutation} of $\{0, 1, \dots, m - 1\}$ called a \emph{probing sequence}:
    \[ h: \mathcal{U} \times \{0, 1, \dots, m - 1\} \to \{0, 1, \dots, m-  1\} ~ \text{ such that } ~ k ~ \text{ has probing sequence } (h(k, 0), h(k, 1), \dots, h(k, m - 1)). \]
    We implement our desired operations as follows:
    \begin{itemize}
        \item $\textsc{Insert}(k, v)$: Find the smallest $p$ such that $T[h(k, p)] = \textsc{Null} \text{ or } \textsc{Deleted}$, then set $T[h(k, p)] = (k, v)$.
        \item $\textsc{Search}(k)$: Find the smallest $p$ such that $T[h(k, p)].\textsc{Key} = k$, then return $T[h(k, p)].\textsc{Value}$.
        \item $\textsc{Delete}(k)$: Find the smallest $p$ such that $T[h(k, p)].\textsc{Key} = k$, then set $T[h(k, p)] = \textsc{Deleted}$.
    \end{itemize}
\end{solution}

\begin{remark}
    Some CYU questions:
    \begin{itemize}
        \item Why do we need a $\textsc{Deleted}$ type?  What goes wrong if we set $T[h(k, p)] = \textsc{Null}$ when deleting?
        \item What happens when $\alpha = \frac{n}{m}$ approaches $1$, or exceeds $1$?
    \end{itemize}
\end{remark}

\begin{claim}
    Suppose that every key receives an \emph{independently generated uniformly random sequence}.
    
    Then the expected number of $p$ we will need to check is $O\left(\frac{1}{1 - \alpha}\right)$, where $\alpha := \frac{n}{m} < 1$.
\end{claim}

\begin{proof}
    Handwaving\dots a guess for $p$ succeeds with probability $1 - \alpha$, so the EV \# of tries is $\frac{1}{1 - \alpha}$. ~ $\blacksquare$

    The rigorous proof goes $\underset{h}{\mathbb{E}}[X] = \displaystyle\sum_{j = 1}^m j \cdot \underset{h}{\mathrm{Pr}}[X = j] = \displaystyle \sum_{j = 1}^m \underset{h}{\mathrm{Pr}}[X \geq j] \leq \displaystyle\sum_{j = 1}^m \alpha^{j - 1} \leq \frac{1}{1 - \alpha}$, hooray. ~ $\blacksquare$

    Even more formally, we need to justify $\underset{h}{\mathrm{Pr}}[X \geq j] \leq \alpha^{j - 1}$, so we should be writing something like the following, where $E_i$ is the event that the $i^{\text{th}}$ probe finds an occupied slot:
    \[ \underset{h}{\mathrm{Pr}}[X \geq j] = \underset{h}{\mathrm{Pr}}[E_1 \wedge E_2 \wedge \dots \wedge E_{j - 1}] = \underset{h}{\mathrm{Pr}}[E_1] \times \underset{h}{\mathrm{Pr}}[E_2 \mid E_1] \times \dots \times \underset{h}{\mathrm{Pr}}[E_{j - 1} \mid E_1 \wedge E_2 \wedge \dots \wedge E_{j - 2}]. \]
    And the $\ell^{\text{th}}$ term in this product is upper-bounded by $\frac{n - \ell + 1}{m - \ell + 1} \leq \frac{n}{m} = \alpha$, as desired. ~ $\blacksquare$
\end{proof}

Generating probing sequences uniformly at random is hard, so in practice, we implement like so.

\begin{solution}{Probing Methods}
    Pick any two hash functions $h_1, h_2: \mathcal{U} \to \{0, 1, \dots, m - 1\}$.  We can use:
    \begin{itemize}
        \item \textbf{Linear Probing.} Take $h(k, p) := (h_1(k) + p) \%m$.
        \item \textbf{Double Hashing Probing.} Take $h(k, p) := (h_1(k) + p \cdot h_2(k)) \%m$.  (This is a permutation only when $\gcd(h_2(k), m) = 1$.)
    \end{itemize}
\end{solution}

% \begin{remark}
%     If $\alpha$ gets too large, we should \emph{double} the value of $m$.  Too small, we \emph{halve} it.  (Recitation \#3) \textcolor{red}{TODO}
% \end{remark}

Now for something new.  Let's try to implement\dots

\begin{problem}{Static Dictionaries \& Perfect Hashing}
    We are given $\{(k_1, v_1), \dots, (k_n, v_n)\}$ up front, after which point only $\textsc{Search}(k)$ may be called.  We would like the following:
    \[ \text{Expected Pre-Processing Time: } O(n) ~~ \mid ~~ \text{Space: } O(n) ~~ \mid ~~ \text{Search Operations: } O(1) \text{ always} \]
    In particular, we require \emph{perfect hashing}: no collisions occur at all!
\end{problem}

Here's a naive solution that passes all constraints except the space requirement.

\begin{solution}{Suboptimal-Space Solution}
    Consider a universal hash family $\mathcal{H}$, and take $m \geq n^2$.

    Then pick $h \in \mathcal{H}$, and use $k_i \mapsto h(k_i)$.  If the first $h \in \mathcal{H}$ fails, pick another $h \sim \mathcal{H}$ until success.
\end{solution}

\begin{proof}
    This works in expected $O(n)$ pre-processing time because $\underset{h \sim \mathcal{H}}{\mathrm{Pr}}[\text{no collisions}] \geq \frac{1}{2}$ by Union Bound---try it yourself!  (When does \emph{universality} of $\mathcal{H}$ get used?) ~ $\blacksquare$
\end{proof}

Here's the smart solution by Fredman-Komlos-Szemeredi (1984) that passes the space requirement.

\begin{solution}{Optimal-Space Solution}
    Sample a hash function $h_1 \in \mathcal{H}_1$ that goes $h_1: \mathcal{U} \to \{0, 1, \dots, n - 1\}$.
    \begin{itemize}
        \item Suppose $N_{\ell} := |h_1^{-1}(\ell) \cap \{k_1, \dots, k_n\}|$ for all $\ell \in \{0, 1, \dots, n - 1\}$.
        \item Use the suboptimal-space solution to fit $N_{\ell}$ key-value pairs in the $\ell^{\text{th}}$ bucket.
        
        Using $O(N_{\ell})$ time and $O(N_{\ell}^2)$ space, give the $\ell^{\text{th}}$ bucket a sampled hash function $h_{2, \ell} \in \mathcal{H}_{2, \ell}$ that goes $h_{2, \ell}: h_1^{-1}(\ell) \to \{0, 1, \dots, N_{\ell}^2 - 1\}$.
    \end{itemize}
    Our data structure is then a \emph{two-dimensional} array: $n$ buckets, the $\ell^{\text{th}}$ of which has $N_{\ell}^2$ slots to hash into.
\end{solution}

\begin{proof}
    It's easy to see that this has $O(1)$ search time, and that we only need $O(n)$ pre-processing time once $h_1$ is chosen.

    The hard part is achieving $O(n)$ space.  If we choose $h_1$ unluckily, we won't actually get $O(n)$ space!  However,
    \[ \textbf{Claim.} ~~~~~ \underset{h \sim \mathcal{H}_1}{\mathrm{Pr}} \left [ \left(\sum_{\ell = 0}^{n - 1} N_{\ell}^2\right) > 4n \right ] < \dfrac{1}{2}. \]
    To prove this, note that by Markov's inequality, it suffices to show $\underset{h \sim \mathcal{H}_1}{\mathbb{E}}\left [ \sum_{\ell = 0}^{n - 1} N_{\ell}^2 \right ] < 2n$.  This takes some clever work:
    \begin{align*}
        \underset{h \sim \mathcal{H}_1}{\mathbb{E}} \left [ \sum_{\ell = 0}^{n - 1} N_{\ell}^2 \right ] = \ & \underset{h \sim \mathcal{H}_1}{\mathbb{E}} \left [ \sum_{\ell = 0}^{n - 1} \left ( N_{\ell} + \sum_{j \neq j'} \mathds{1}\left [ h(k_j) = h(k_{j'}) = \ell \right ] \right ) \right ] \\
        = \ & \underset{h \sim \mathcal{H}_1}{\mathbb{E}}\left [ \sum_{\ell = 0}^{n - 1} N_{\ell} \right ] + \underset{h \sim \mathcal{H}_1}{\mathbb{E}} \left [  \sum_{\ell = 0}^{n - 1}  \sum_{j \neq j'} \mathds{1}\left [ h(k_j) = h(k_{j'}) = \ell \right ] \right ] \\
        = \ & n + \sum_{j \neq j'} \left(\underset{h \sim \mathcal{H}_1}{\mathrm{Pr}}[h(k_j) = h(k_{j'})] \right) \\
        \leq \ & n + \frac{n(n - 1)}{n} < 2n.
    \end{align*}
    Thus, we can guarantee $O(n)$ space in expected $O(n)$ time by repeatedly resampling $h_1$ until the space is less than $4n$. ~ $\blacksquare$
\end{proof}


\end{document}